% Parts A and E: working text. Root supplies macros/preamble.
\section{Strict bifibrations and equivalence transports}
\subsection{The two-type obstruction}

\begin{result}{A.1}{Proved}{Equivalence transports have no higher connecting maps}
Let $B$ and $C$ be small categories in the fixed universe. Let $p:E\to B$ be a Grothendieck opfibration whose fibre categories are equivalent to $C$ and whose transport functors are equivalences of categories. For every basepoint $b$ and every object $x$ of its fibre, the connecting homomorphisms
\[
 \partial:\pi_{r+1}(BB,b)\longrightarrow\pi_r(BC,x)
\]
of its classifying-space homotopy-fibre sequence are zero for $r\geq2$. Consequently, if $BE$ is weakly contractible and $B$ is nonempty and connected, then $BC$ is a $1$-type and $BB$ is a $2$-type. In particular there is no such opfibration with weakly contractible total category over the five-element rectangular-band monoid of the question, whose classifying space is $S^2$.
\end{result}
\begin{proof}
Let $\operatorname{Eq}(C)$ be the monoidal groupoid of equivalences $C\to C$ and natural isomorphisms. Its monoidal nerve is a group-like space. Its delooping $B\operatorname{Eq}(C)$ is a $2$-type: its loop space is the nerve of a groupoid and therefore a $1$-type. A cleavage and fibre identifications give a pseudofunctor from $B$ to the one-object bicategory with endomorphism groupoid $\operatorname{Eq}(C)$. The associativity and unit isomorphisms are precisely the coherence needed to define its classifying map $BB\to B\operatorname{Eq}(C)$. Applying classifying spaces to equivalences and natural isomorphisms gives a map
\[
 B\operatorname{Eq}(C)\longrightarrow B\operatorname{hAut}(BC).
\]
The local system underlying $BE\to BB$ is the pullback of the universal family along the composite. This follows directly from the Grothendieck homotopy-colimit description: both families have fibre $BC$, the prescribed transports, and the same coherent transition maps. Naturality of the homotopy-fibre sequence factors the displayed connecting map through $\pi_{r+1}(B\operatorname{Eq}(C))$, which vanishes for $r+1\geq3$.

When $BE$ is weakly contractible, its long exact sequence makes every displayed connecting map an isomorphism for $r\geq2$. Hence both its source and target vanish. This proves the two truncation assertions. For $BB\simeq S^2$, $\pi_3(BB)\cong\mathbb Z$, a contradiction. Empty fibres are irrelevant to a contractible nonempty total space over a connected nonempty base: equivalence transport and zigzags in the base make either all fibres empty or all fibres nonempty. For a disconnected base the argument applies separately on each component with a nonempty fibre.
\end{proof}

\begin{result}{A.2}{Proved}{All steps of the central-cocycle argument in Lead 2}
Let $M=\{1,e_{00},e_{01},e_{10},e_{11}\}$ with $e_{ij}e_{kl}=e_{il}$. If the transport functors of a Grothendieck opfibration $E\to\mathbf BM$ are categorical equivalences, there is a normalized central cocycle $c:M\times M\to A$, where $A=\operatorname{Aut}(1_C)$, and an isomorphism over $\mathbf BM$ from $E$ to the category having objects those of $C$ and arrows $(m,f):x\to y$, $f:x\to y$ in $C$, with composition twisted by $c$. The action in Lead 2 is a degreewise free action of the simplicial group $N(\mathbf BA)$; with that precise meaning the Borel argument and its conclusion are valid.
\end{result}
\begin{proof}
Choose identity cocartesian lifts over $1$. Each nonidentity $m$ is idempotent, so its coherence is an isomorphism $T_mT_m\cong T_m$. Since $T_m$ is an equivalence, cancellation gives a natural isomorphism $T_m\cong1_C$. Choose such isomorphisms and transfer the pseudofunctor structure to the identity functors. A natural automorphism of $1_C$ commutes with every arrow of $C$, and two such automorphisms commute componentwise by naturality. Thus $A$ is abelian. In additive notation the transferred associativity equation is
\[
 c(k,m)+c(km,n)=c(m,n)+c(k,mn),\qquad c(1,m)=c(m,1)=0.
\]
Every arrow of $E$ over $m$ factors uniquely through the chosen cocartesian arrow; the chosen isomorphism $T_mx\cong x$ then identifies it uniquely with an arrow $f:x\to y$ of $C$. This is an identity-on-objects bijection on each hom-set. With the convention that the later arrow multiplies on the left, its composition law is
\[
 (n,g)\circ(m,f)=(nm,gf\,c(n,m)_x),
\]
where the central automorphism is written multiplicatively when acting on $x$. It is therefore a categorical isomorphism, not merely an equivalence, after the choices have been made.

Let $M_c=M\times A$ with multiplication
$(m,a)(n,b)=(mn,a+b+c(m,n))$, using the same composition-order convention throughout. The functor $C\times\mathbf B M_c\to E$ sends $(f,(m,a))$ to $(m,f a)$. The simplicial group $G=N(\mathbf BA)$ acts on $N(C\times\mathbf B M_c)$: in degree $q$, the element $(b_1,\ldots,b_q)\in A^q$ multiplies the $q$ successive $C$-arrows by $b_i$ and subtracts $b_i$ from their $M_c$-coordinates. Centrality makes this commute with composition faces; normalization makes it commute with degeneracies. The action is free because the $A$-coordinates in $M_c$ recover every $b_i$. Its orbit simplicial set is exactly $NE$, degree by degree.

To justify the homotopy quotient without an unstated local-triviality hypothesis, form the bisimplicial bar construction for this action. In every fixed simplicial degree its action groupoid is a free action groupoid, and its nerve maps by a weak equivalence to the discrete orbit set. The diagonal weak-equivalence theorem for bisimplicial sets gives
\[
 B E\simeq \bigl(BC\times B M_c\bigr)\mathbin{//}BA.
\]
The homotopy-quotient fibre sequence is
$BC\times BM_c\to BE\to B^2A$. If $BE$ is contractible, its fibre is weakly equivalent to $\Omega B^2A=BA$. Thus $\pi_r(BC\times BM_c)=0$ for $r\geq2$, and projection onto the first factor makes $\pi_r(BC)=0$. This verifies the claimed implication. An ordinary discrete $A$-action would be the wrong statement; the simplicial group $N(\mathbf BA)$ is essential.
\end{proof}

\begin{result}{A.3}{Proved}{Exact classification for equivalence transports}
Fix a connected nonempty small category $B$ and a locally constant space-valued object with fibre homotopy type $X$. It has a bifibration representative whose transports are equivalences of categories if and only if there are a small category $C$, an equivalence $BC\simeq X$ of spaces, and a homotopy lift of its classifying map
\[
 BB\longrightarrow B\operatorname{hAut}(X)
\]
through $B\operatorname{Eq}(C)\to B\operatorname{hAut}(BC)\simeq B\operatorname{hAut}(X)$. For disconnected $B$, this condition is applied independently to every component, allowing a different $C$ on different components.
\end{result}
\begin{proof}
Necessity is the classifying-map construction of A.1, with the chosen fibre comparison retained. Conversely the nerve of the one-object bigroupoid $\operatorname{Eq}(C)$ is a Kan $2$-type. A homotopy class of maps from $NB$ into it is represented by an actual simplicial map, because $NB$ is cofibrant and the target is Kan. Such a map is precisely a normal pseudofunctor: its edges are equivalences $C\to C$, its triangles are their invertible composition constraints, and its tetrahedra assert the coherence equation. The higher simplices are determined by these data. Its Grothendieck construction is an opfibration. Every equivalence has a right adjoint, so the same projection is a fibration. The given homotopy between classifying maps supplies the required equivalence with the specified local system. For $X=\varnothing$, take $C=\varnothing$; the same argument applies. The condition is a lift through the indicated categorical automorphism $2$-group, not mere factorization through some unspecified $2$-type.
\end{proof}

\subsection{Splitting the rectangular-band action}
\begin{result}{A.4}{Proved}{The split data and their curvature class}
Up to equivalence of coherent pseudo-actions, an action of $M$ on a small category $C$ is specified by a small category $D$, functors $A_i:D\to C$, $B_j:C\to D$, and natural isomorphisms $\beta_{ji}:B_jA_i\to1_D$, for $i,j\in\{0,1\}$. The corresponding nonunit functors are $T_{ij}=A_iB_j$ and their multiplication constraints are
\[
 (A_iB_j)(A_kB_l)\xrightarrow{A_i\beta_{jk}B_l}A_iB_l.
\]
There is no additional equation on the four $\beta_{ji}$: the pseudo-action associativity equation follows by interchange for two disjoint contractions. Equality $T_{ij}=A_iB_j$ refers to this equivalent model, not to the original unmodified functors.
\end{result}
\begin{proof}
We give the splitting used in the comparison. Write $P=T_{00}$ and $\mu:P^2\to P$ for its multiplication isomorphism. Associativity and cancellation of $\mu$ imply $P\mu=\mu P$. Let $D$ have objects $(x,\theta)$ with $\theta:Px\to x$ an isomorphism satisfying $P\theta=\mu_x$, and arrows $h:(x,\theta)\to(y,\psi)$ satisfying $h\theta=\psi Ph$. Let $A:D\to C$ forget $\theta$, and let $B:C\to D$ send $x$ to $(Px,\mu_x)$. Then $AB=P$ strictly. The arrows $\theta$ define an isomorphism $\epsilon:BA\to1_D$, and $A\epsilon B=\mu$. Every category and functor just defined is small.

Equivalently, this extends the action to the two-object subcategory of the idempotent completion of $\mathbf BM$ on $1$ and $e_{00}$. Its arrows from $e_{00}$ to $1$ are $e_{i0}$, and those from $1$ to $e_{00}$ are $e_{0j}$. Their composites satisfy $e_{0j}e_{i0}=e_{00}$, the identity of the new object. Explicit choices are $A_i=T_{i0}A$ and $B_j=BT_{0j}$, with the isomorphisms to $T_{ij}$ and to $1_D$ obtained from the original multiplication constraints and $\epsilon$. All their coherence diagrams are images of the original associativity diagrams together with $A\epsilon B=\mu$. Conversely the displayed formula defines a normal pseudo-action from the split data: on $A_iB_jA_kB_lA_rB_s$, contracting $B_jA_k$ and $B_lA_r$ in either order gives the same natural isomorphism by naturality and interchange. The split construction and this reconstruction are inverse up to the specified coherent equivalences; a different splitting is related by the functors induced from $B'A$ and $BA'$ and their counit isomorphisms.
\end{proof}

For these data, the exact adjoint condition is representability of
\[
 x\longmapsto\operatorname{Hom}_C(A_iB_jx,y)\qquad(C^{\mathrm{op}}\to\mathbf{Set})
\]
for every $i,j$ and every $y\in C$. Representing objects $G_{ij}y$ assemble uniquely to right adjoints by Yoneda. If these right adjoints exist, $B(A_iB_j)$ is a homotopy equivalence by its adjunction. Since $B(B_jA_i)\simeq1$, each $BA_i$ and $BB_j$ is a homotopy equivalence. Put $X=BD$ and write $a_i,b_j$ for these space maps.

The four isomorphisms determine the following precise clutching element. Choose a homotopy $\eta:a_0b_0\simeq1_{BC}$ completing $(a_0,b_0,\beta_{00})$ to a homotopy equivalence; choose it with its triangle coherence. Define the path of equivalences
\[
 u_1:a_1\xRightarrow{\eta^{-1}a_1}a_0b_0a_1
                \xRightarrow{a_0\beta_{01}}a_0.
\]
Then, in traversal order, the loop
\[
 1_X\xRightarrow{\beta_{10}^{-1}}b_1a_0
 \xRightarrow{(b_1u_1)^{-1}}b_1a_1
 \xRightarrow{\beta_{11}}1_X                                      \tag{A.4.1}
\]
represents the clutching element in $\pi_1\operatorname{hAut}(X)$, with the sphere orientation fixed by this formula. Indeed the two-object category above has four elementary relations $B_jA_i=1_D$. Contracting its $A_0$ edge and using the first three relations as the reference trivializations leaves precisely the fourth relation as the boundary loop (A.4.1). The corresponding normalized bar $2$-cycle in $BM$ is
\[
 [e_{00}|e_{00}]-[e_{00}|e_{10}]-[e_{01}|e_{00}]+[e_{01}|e_{10}].
\]
In the special case where all four functors are identities, evaluation is $\beta_{00}-\beta_{01}-\beta_{10}+\beta_{11}$. This fixes the sign. Changing the fibre framing conjugates the clutching element by the chosen equivalence of $X$; changing the sphere orientation inverts it.

The exact contractibility condition is also expressible in this language. Let $\chi:S^2\to B\operatorname{hAut}(X)$ be the class represented by (A.4.1), and choose $x\in X$. The orbit map
\[
 \Omega S^2\xrightarrow{\Omega\chi}\operatorname{hAut}(X)
                  \xrightarrow{\operatorname{ev}_x}X             \tag{A.4.2}
\]
must be a weak equivalence. This condition is necessary and sufficient: choosing a point in the associated total space gives a map from the based path fibration over $S^2$ into that family, whose fibre map is (A.4.2); the fibrewise equivalence criterion proves both directions. The path local system has $X=\Omega S^2$ and clutching loop given by left translation along a generator of $\pi_1\Omega S^2$. Merely requiring evaluation of (A.4.1) to generate $\pi_1X$ is weaker than (A.4.2) and is not asserted sufficient.

\subsection{The complete monoid enumeration through order five}
\begin{result}{A.5}{Proved}{The unique non-$1$-type monoid of order at most five}
Up to monoid isomorphism, the rectangular-band monoid $M$ is the only monoid with at most five elements whose classifying space is not a $1$-type. The numbers of isomorphism classes in orders $1,2,3,4,5$ are respectively $1,2,7,35,228$; the complete multiplication tables and reduction certificates accompany the paper.
\end{result}
\begin{proof}
We first justify the reductions used by the computation. If $e$ is an idempotent in a monoid $L$ and $eLe=eL$ or $eLe=Le$, then $BL\simeq B(eLe)$. Form the two-object category $K$ with objects $1,e$ and arrows $K(x,y)=yLx$. The inclusion of the object $1$ is homotopy initial: its comma category over $1$ has a terminal object, and the category over $e$ has objects $m\in eL$ and arrows $a:m\to m'$ with $m'a=m$. The components $m:m\to e$ give a natural transformation from its identity functor to the constant functor at $e$, so that comma category is contractible. If $eLe=Le$, the inclusion of the object $e$ is homotopy initial: in its comma category over $1$, the object $e$ is terminal, because every $m\in Le$ belongs to $eLe$ and factors uniquely through $e$. If $eLe=eL$, use instead the outgoing comma category, in which $e$ is initial. Quillen's Theorem A, and its opposite, give the claimed zigzag through $BK$. For $e\ne1$, the corner has strictly fewer elements, since $1\in eLe$ would imply $e=1$.

The program fixes the identity as $0$ and enumerates every multiplication table. At each node it assigns an unknown entry and repeatedly imposes all currently determined associativity equalities $(ab)c=a(bc)$; a conflict rejects that branch. No value is discarded except by one of these equalities. Hence every associative table occurs. The terminal tables are checked again against all $n^3$ equations and quotiented by all permutations fixing $0$. The labeled counts are $1,2,11,156,4122$. The independent verification program then records every permitted corner reduction and checks its target against the same table catalogue. There remain exactly the trivial group, $C_2,C_3,C_4,C_2\times C_2,C_5$, and $M$. Of the $273$ total isomorphism classes, precisely one reduces to $M$, namely $M$ itself. All other types reduce to groups and are $1$-types.

For completeness, $\pi_1(BM)$ is trivial because every nonunit generator is idempotent. Put $M_i=\{e_{i0},e_{i1}\}$. The right $\mathbb Z[M]$-modules $\mathbb Z[M_i]$ are projective idempotent images. The exact resolution
\[
0\to\mathbb Z[M_0]\oplus\mathbb Z[M_1]\to\mathbb Z[M]
 \xrightarrow{e_{00}-e_{01}}\mathbb Z[M_0]\to\mathbb Z\to0
\]
has the first map the two inclusions and the last the augmentation. The middle map sends $1$ to $e_{00}-e_{01}$ and all nonunits to zero, proving exactness directly. Tensoring with the trivial left module gives $\mathbb Z^2\xrightarrow{(1,1)}\mathbb Z\xrightarrow0\mathbb Z$. Thus $BM$ has the integral homology of $S^2$; simple connectivity and the homological Whitehead theorem give $BM\simeq S^2$. This is also the explicit unnumbered lemma in Fiedorowicz (2002, pp.~33--34), whose primary text was opened and checked.
\end{proof}

\begin{result}{A.6}{Known}{Every local system over a $1$-type base has a strict bifibration model}
If $BB$ is a $1$-type, every small space-valued locally constant object over $B$ has a small strict bifibration representative with categorical-isomorphism transports. This includes empty values and disconnected bases.
\end{result}
\begin{proof}
Let $q:B\to\Pi_1B$ be ordinary groupoid completion. The specified map $Bq$ is a weak equivalence under the hypothesis, so restriction of local systems along it is an equivalence. On each component of $\Pi_1B$ choose a group model $G$ and a Kan fibration of small simplicial sets $Y\to NG$ representing the resulting family. Pull back along the standard contractible free $G$-simplicial set $EG\to NG$. The resulting simplicial set $\widetilde Y$ has a strict $G$-action, is weakly equivalent to the fibre, and its Borel family recovers $Y\to NG$. The simplex category $\Delta/\widetilde Y$ inherits a strict action by categorical isomorphisms. Its classifying space is naturally weakly equivalent to $\widetilde Y$. Integrating this action and pulling back along $q$ gives the required bifibration and its specified diagramwise comparison. A small family of small simplicial sets gives small simplex categories in the fixed universe. Empty fibres yield empty simplex categories. This is the result of STRATA E.9, here proved using simplex categories in place of contractible-patch categories.
\end{proof}

\begin{result}{A.7}{Proved}{Finite acyclic fibre categories cannot solve the sphere test}
Here a finite acyclic category means a finite category with no nonidentity endomorphisms and no directed cycles of nonidentity arrows. A bifibration over $M$ with weakly contractible total category cannot have such a category as fibre; in particular its fibre cannot be a finite poset.
\end{result}
\begin{proof}
Adjoint transport functors are homotopy equivalences, so the Grothendieck homotopy-fibre theorem identifies the fibre with $\Omega S^2$. The homology Serre spectral sequence of $\Omega S^2\to PS^2\to S^2$ has only columns $0$ and $2$ and contractible total space. Its $d_2$ therefore identifies $H_q(\Omega S^2;\mathbb Z)$ with $H_{q+1}(\Omega S^2;\mathbb Z)$ for every $q\geq0$. Since $\Omega S^2$ is connected, these groups are all $\mathbb Z$. A finite acyclic category has a finite-dimensional nerve: a nondegenerate simplex cannot revisit an object. Its integral homology vanishes above its number of objects minus one, a contradiction.
\end{proof}

\begin{result}{A.8}{Open}{The remaining strict-adjoint construction and the requested action search}
We do not construct, or exclude, a bifibration with weakly contractible total category over $M$ when its transports have right adjoints but are not equivalences. The same question for the seven-element monoid of CATHOT G.10 remains open for such transports; A.1 excludes equivalence transports there because $\pi_3(S^3)\ne0$. The general-base existence question and the exhaustive pseudo-action enumeration on categories with at most five objects and twelve arrows remain unresolved here. In the proposed filtered-colimit enlargement of the fence category, a definition of the completion together with endpoint-compatible universal arrows and a fibrewise weak-equivalence comparison is still missing. Closing under all presheaves would produce contractible fibre categories and would therefore lose the required fibre $\Omega S^2$; the smaller filtered completion is not proved to acquire the adjoints. For categories of reduced words, no specified composition and adjoint pair satisfying all rectangular-band relations is constructed here. For a bar or bisimplicial resolution, its contraction does not by itself give both categorical lifting properties of the endpoint projection; no shift construction with both universal properties is established. The four Beck--Chevalley squares have not been proved invertible from the adjunction data. Accordingly none of these routes is recorded as a negative existence theorem: each still lacks its decisive comparison or adjunction step. A.5 completes only the requested enumeration of base monoids of order at most five. No finite pseudo-action table is asserted to be exhaustive.
\end{result}

\section{The full path categories}
\subsection{The three-object marked chain}
\begin{result}{E.1}{Proved}{All alternating fixed-shape path spaces for the marked chain}
Let $P=[2]$ and let $W$ consist of its identities and $0\to2$. For $n\geq2$, let $F_n(a,b)$ be the full poset of paths of type $(-,+)^n$ from $a$ to $b$, with backward arrows in $W$ and arbitrary endpoint-fixed vertical arrows. Every $F_n(a,b)$ is contractible, except $F_n(1,1)$, which is the disjoint union of an isolated point and a contractible nonempty poset. These two components are preserved by identity padding. The corresponding assertions hold integrally in every homological degree and for all homotopy groups, not only through a finite range.
\end{result}
\begin{proof}
An object is a sequence $(x_0,\ldots,x_{2n})$ with $x_0=a$, $x_{2n}=b$. At every negative edge the pair $(x_{2j},x_{2j+1})$ is one of $(0,0),(1,1),(2,2),(2,0)$; at every positive edge $x_{2j+1}\leq x_{2j+2}$. Morphisms are pointwise inequalities.

If $a=0$, the path $(0,0,\ldots,0,b)$ is least. If $a=2$, the path $(2,0,\ldots,0,b)$ is least. For $a=1,b=2$, the path $(1,1,2,\ldots,2)$ is greatest. For $a=1,b=0$ and $n\geq2$, the path $(1,1,2,\ldots,2,0,0)$ is greatest. Each satisfies all edge conditions, and comparison coordinate by coordinate proves its stated universal property.

It remains to treat $(a,b)=(1,1)$. Any nonconstant path must contain both a $2$ and a $0$: it can leave $1$ only by a positive move to $2$, and it can return below $2$ only through the allowed negative move $2\to0$. Such a path is incomparable with the constant path, which is therefore an isolated component. Let $Q_n$ be the remaining full subposet. If $p\in Q_n$, its first zero coordinate $k(p)$ is odd and at least $3$. Define $R(p)$ by replacing coordinates $2,\ldots,k(p)-1$ by $2$ and leaving every other coordinate unchanged. This is an admissible path, since its modified prefix is $1,1,2,\ldots,2,0$, followed by the original suffix.

We have $p\leq R(p)$. If $p\leq q$, then $k(p)\leq k(q)$, because every zero of $q$ is a zero of $p$. Before $k(p)$ both modified paths have the same maximal interior value. Between $k(p)$ and $k(q)$, the coordinates of $R(q)$ are $2$; from $k(q)$ onward the comparison is the original comparison $p\leq q$. Hence $R$ is monotone. Its image is exactly the principal upper set of
\[
 u_n=(1,1,2,0,\ldots,0,1).
\]
Indeed a path above $u_n$ has coordinate $2$ equal to $2$, and must remain $2$ until its first zero; it is therefore fixed by $R$. Conversely every $R(p)$ lies above $u_n$. Thus $R$ is a closure operator with image having initial object $u_n$, and $NQ_n$ is contractible. Padding a path by a final identity pair preserves whether it is constant. It therefore preserves the two displayed components. This proves all claims, including the eventual stability in $n$.
\end{proof}

\begin{result}{E.2}{Proved}{The comparison $J$ is an equivalence for the marked chain}
For $([2],\{\mathrm{id},0\to2\})$, the specified comparison from weak-vertical paths to full paths is a weak equivalence for every pair of endpoints. Its infinity-localization is equivalent to the ordinary walking-retraction category, with objects $A,B$, arrows $i:A\to B$, $r:B\to A$, and $ri=1_A$. Its mapping spaces are discrete; $\operatorname{End}(B)=\{1_B,ir\}$, and every other hom-set is a singleton. The objects $0,2$ map to $A$ and $1$ maps to $B$.
\end{result}
\begin{proof}
We first specify the passage from E.1 to the all-type category. Let $\mathrm{Ty}$ have sign words as objects and monotone endpoint-preserving type maps as arrows. Enrich each hom-set to the poset of such maps and endpoint-fixed transformations. We check its contraction explicitly. If the target word is empty, this poset is terminal. If the first source run has sign opposite to the first target run, every type map is constant on that source run, so deleting it is an isomorphism of hom-posets. Otherwise let the first source and target runs have vertex intervals $[0,a]$ and $[0,b]$ with common sign $\varepsilon$. The map $h(j)=\max(j,b)$ collapses the first target run and fixes the remaining vertices. For a type map $F$, put $RF(0)=0$ and $RF(i)=hF(i)$ for $i>0$. No map crosses the first target turning vertex during its first source run. Consequently $RF$ has that first run equal to $(0,b,\ldots,b)$; it is a monotone endpoint-preserving functor. Away from vertex $0$, apply $h$ to transformations. If $\varepsilon=+$, the pointwise arrows give $1\Rightarrow R$; if $\varepsilon=-$, they give $R\Rightarrow1$. At both endpoints these components are identities, and on the first source edge the required square commutes because the target is thin. The image is the full subposet with the prescribed first run, on which $R$ is the identity. Deleting the two prescribed runs reduces the number of target runs and completes the induction. Each type map preserves backward edges as a string of backward edges or an identity; thus this argument does not introduce unmarked backward edges. Precomposition of full path categories is an action of these hom-posets; arbitrary vertical components are exactly what makes this enriched action defined. Passing to classifying spaces therefore extends the diagram coherently to the reachability preorder of types, since the enriched mapping spaces are empty or contractible.

The functor from $\mathrm{Ty}^{\mathrm{op}}$ to this preorder is homotopy final. At a type $\sigma$, its comma category is the opposite of the full subcategory of types admitting a map to $\sigma$. The functor $\tau\mapsto\sigma*\tau$ has transformations to $\tau$ and to $\sigma$ by block collapse, so this comma category is contractible. The alternating sequence $t_n=(-,+)^n$ is also final in the preorder: every finite sign word is obtained from some $t_n$ by collapsing unused edges and subdividing retained runs, and the permissible values of $n$ form a nonempty final segment of $\mathbb N$. Thus the canonical Grothendieck homotopy-colimit comparison gives
\[
 B\operatorname{Ch}_W(P;a,b)\simeq\operatorname*{hocolim}_n BF_n(a,b).
\]
This is the enriched-type argument of PATHS C.P2--C.P3, with its comma contractions and the necessary arbitrary-vertical action retained. E.1 shows that the right side is homotopy discrete, with two components exactly at $(a,b)=(1,1)$.

We recall why the specified map $J$ has a retraction on homotopy types, so that no arbitrary comparison replaces it. Let $\mathcal L$ be the infinity-localization. For any type $\tau$, invert its backward edges; its interval becomes the ordinal whose length is the number of forward edges. The universal localization property identifies the corresponding relative diagram category in $\mathcal L$ with the ordinal diagram category. Endpoint composition then gives a functor to $\operatorname{Map}_{\mathcal L}(a,b)$. These equivalences are natural in type maps, so their inverses can be taken as one inverse equivalence of diagrams. Groupoid completion and the type homotopy colimit give a map
$Q:B\operatorname{Ch}_W(P;a,b)\to\operatorname{Map}_{\mathcal L}(a,b)$.
On weak-vertical paths $QJ$ is the canonical hammock comparison, a weak equivalence. The restricted comparison is Dwyer--Kan's hammock mapping-space construction; its Grothendieck form follows from the Reedy-cofibrant diagram of fixed-shape nerves, whose latching inclusions insert identity edges. This is PATHS B.2 and B.6 with the actual comparison maps retained.

Consequently $J$ is split injective on every higher homotopy group. Since its target has no higher homotopy, neither does its source. Both component sets are the morphisms of the ordinary localization: composing paths after formally inverting $W$ gives that identification, and elementary composition, identity insertion, and cancellation of a marked backtrack give exactly its relations. Thus $J$ is also bijective on components and is a weak equivalence. An infinity-category with discrete mapping spaces is equivalent to the nerve of its homotopy category. Identifying $0$ with $2$ using the marked composite leaves exactly $i,r$ with $ri=1$, which is the displayed walking retraction.
\end{proof}

The program \texttt{code/AE/path\_chain.py} constructs these posets independently for lengths $2,4,6,8,10,12$, checks all pointwise comparisons, and records elementary beat-point-removal certificates. At length $6$, the object counts, with row $a$ and column $b$, are
\[
 \begin{pmatrix}6&9&13\\4&6&9\\9&13&19\end{pmatrix},
\]
and the component counts are $1$ in every entry except $(1,1)$, which has $2$. The proof, rather than numerical stabilization, establishes the unbounded limit. These are the cofinal fixed-shape categories, not the full subcategory containing every word of length at most six; the latter finite computation is not claimed here.

\subsection{The finite catalogue and its limits}
\begin{result}{E.3}{Proved}{Complete relative-poset catalogue through four objects}
There are respectively $1,3,14,102$ isomorphism classes of pairs $(P,W)$ on $1,2,3,4$ objects, where $P$ is a poset and $W$ contains identities and satisfies two-out-of-three. The attached catalogue lists every pair and also the unique empty relative poset, for $121$ types in all. The program computes all endpoint full-path posets of types $(-,+)^n$, $n=1,2,3$, giving $5313$ rows. In $5307$ rows all connected components are certified contractible by beat-point removal. Exactly six rows have a non-discrete remaining core, all for $P=[3]$ and $W=\{\mathrm{id},0\to2,1\to3\}$. These are bounded calculations, not a proof of the all-type comparison for all $120$ pairs.
\end{result}
\begin{proof}
Every finite poset has a linear extension. The program therefore enumerates each subset of pairs $(i,j)$ with $i<j$, retains exactly the transitive relations, and identifies relations under every permutation of the object set. It then enumerates every subset of the nonidentity arrows and retains exactly those for which no composable triple has exactly two of its three arrows marked. This is precisely the two-out-of-three condition, including composition closure. A second canonicalization under all permutations identifies the relative pairs. Thus the finite catalogue is exhaustive.

For every endpoint pair and every stated type, a recursive traversal lists exactly the legal vertex sequences. Pointwise order gives exactly its arrows, because a poset has at most one arrow of each type and every naturality square commutes automatically. Each deleted beat point has either a greatest strict predecessor or a least strict successor; sending it to that object and fixing all others is a deformation retraction of nerves. The recorded sequence of these deletions therefore certifies the stated contractible components. The six exceptional cores are retained in full in the output. Four are the four-point circle. The other two have $10$ vertices, $25$ edges, and $16$ triangles. Exact integral Smith forms give vanishing reduced homology, and the spanning-tree fundamental-group presentation has $16$ generators which the recorded Tietze eliminations remove completely. They are therefore simply connected and acyclic, hence contractible by Hurewicz and Whitehead. The final table \texttt{relative\_posets\_homotopy\_table.csv} records integral homology through degree three and the fundamental group at every basepoint for all $5313$ rows. Stopping the beat-point algorithm alone was not used to conclude this.
\end{proof}

\begin{result}{E.4}{Proved}{Transient higher homotopy in a four-object calculation}
For $P=[3]$ marked only at identities, $0\to2$, and $1\to3$, the full fixed-shape path poset $F_2(1,2)$ of type $(-,+)^2$ has classifying space $S^1$, whereas the infinity-localization of $(P,W)$ is the terminal infinity-category. Thus a fixed finite path shape need not already give the localization mapping space, even in this small example.
\end{result}
\begin{proof}
The path poset has exactly four objects
\[
 (1,1,1,1,2),\quad(1,1,3,1,2),\quad
 (1,1,2,0,2),\quad(1,1,2,2,2).
\]
The first and third are minimal and each is below both of the other two; there are no other nonidentity comparisons. Its nerve is the four-edge circle. Write the successive generating arrows of $[3]$ as $f,g,h$. In the localization, $gf$ and $hg$ are equivalences. Hence $g$ has both a left and a right homotopy inverse and is an equivalence; then so are $f$ and $h$. The localization is therefore the groupoid completion of $[3]$, which is contractible. The program verifies the four objects, all four comparisons, and the exact finite-poset description.
\end{proof}

\begin{result}{E.5}{Refuted}{No universal finite length bound for all four-object cases}
The full path category for every relative poset on at most four objects is not computed by a finite all-word length cutoff, even at the level of components.
\end{result}
\begin{proof}
Take the four-object poset with two minimal objects $a_0,a_1$, two maximal objects $b_0,b_1$, and all four relations $a_i<b_j$, and mark every arrow. Its nerve is a four-edge circle. Since all arrows are marked, the path comparison is the identity and its endomorphism space at any object is $\Omega S^1$, with components indexed by $\mathbb Z$. At any finite word-length cutoff, there are finitely many path objects, so its nerve has finitely many components. Its inclusion into the full path category cannot be a weak equivalence. The order relation and its circle nerve are explicitly checked by the finite catalogue. This does not obstruct computing a limit by a sequence of cutoffs; it refutes treating one such cutoff as a proof of that limit.
\end{proof}

\begin{result}{E.6}{Open}{The general comparison and the additive and pointed refinements}
The specified map $J$ for arbitrary two-out-of-three relative categories, its all-endpoint assertion for the remaining four-object relative-poset cases, and the literal all-vertical and ind-stage transfers in PATHS C.9--C.10 remain unresolved here. E.7--E.8 below give the additive and pointed refinements of the corrected localized hom-spaces. E.2 settles the particular three-object example that PATHS B.8 singled out. E.3 does not settle the unbounded four-object assertion. The standard weak-vertical definition computes the localization and has strict concatenation, but a general additive or smash-product refinement requires additional coherent structure; none is deduced solely from the computations above.
\end{result}

\subsection{Additive and pointed refinements of the corrected model}
\begin{result}{E.7}{Proved}{The corrected additive refinement under direct-sum stability}
Let $k$ be a commutative ring in the fixed universe, let $C$ be a small additive $k$-linear category, and let $W$ be a wide subcategory closed under finite direct sums of arrows. The infinity-localization $\mathcal L=C[W^{-1}]$ has a coherent enrichment in the connective derived infinity-category of $k$-modules, written $D^{\leq0}(k)$ in cohomological grading. Its underlying mapping spaces are the localization mapping spaces, and its composition is $k$-bilinear. Consequently the corrected weak-vertical path categories admit composition-compatible $D^{\leq0}(k)$-refinements whose underlying spaces are their nerves. No model-category or suspension--loop hypothesis is required for this assertion about the corrected model.
\end{result}
\begin{proof}
We first check precisely which algebraic structure survives localization. Localization commutes with finite products of relative categories. Indeed a functor $C\times D\to E$ inverts $(W,1)$ and $(1,V)$ exactly when its curried functor takes values in the full subcategory $\operatorname{Fun}_V(D,E)$ and inverts $W$ there. Applying the universal property twice gives the product assertion, including natural transformations and their higher coherences.

The direct-sum functor $\oplus:C\times C\to C$ is relative by the hypothesis, and the diagonal functor is relative. Its ordinary adjunctions $\oplus\dashv\Delta\dashv\oplus$ descend to the localizations: the functors descend, and full faithfulness of restriction along a localization extends their units, counits, and triangle identities. Similarly the two adjunctions defining the zero object descend. Hence $\mathcal L$ has the finite biproducts induced from $C$.

Let $\mathcal F_k$ be the small category of finite free $k$-modules with all matrices, monoidal under tensor product. The ordinary $k$-linear category $C$ has an action
\[
 \mathcal F_k\times C\longrightarrow C,\qquad (k^n,a)\longmapsto a^{\oplus n}.
\]
A matrix acts by its scalar multiples of identity arrows. Matrix multiplication and the distributivity isomorphisms supply its associative and unital module coherence. For every $V\in\mathcal F_k$, the functor $V\otimes-$ preserves $W$, by finite-direct-sum stability. The product-localization argument therefore gives a coherent action of $\mathcal F_k$ on $\mathcal L$, preserving finite direct sums in both variables.

For $a,b\in\mathcal L$ define the presheaf
\[
 H(a,b):\mathcal F_k^{\mathrm{op}}\longrightarrow\mathcal S,
 \qquad H(a,b)(V)=\operatorname{Map}_{\mathcal L}(V\otimes a,b).
\]
It takes finite direct sums in $\mathcal F_k$ to products, because the action takes them to coproducts in $\mathcal L$. The infinity-category of these product-preserving presheaves is $D^{\leq0}(k)$. To see the precise input, $\mathcal F_k^{\mathrm{op}}$ is the algebraic theory of $k$-modules. Its strict simplicial algebras are simplicial $k$-modules. Badzioch's rigidification theorem, Theorems~1.3 and~6.4, identifies homotopy product-preserving diagrams with those strict algebras up to weak equivalence; Dold--Kan identifies the latter with connective chain complexes up to quasi-isomorphism. Thus $H(a,b)$ is an actual object of the claimed derived infinity-category, not only an abelian group structure on $\pi_0$. Evaluation at $k$ gives exactly $\operatorname{Map}_{\mathcal L}(a,b)$.

We also construct composition, since objectwise rigidification alone would not suffice. For $U,V\in\mathcal F_k$, composition and the action give the natural pairing
\[
 \operatorname{Map}(U\otimes b,c)\times\operatorname{Map}(V\otimes a,b)
 \longrightarrow\operatorname{Map}((U\otimes_kV)\otimes a,c),
 \quad(f,g)\longmapsto f\circ(U\otimes g).
\]
This defines, by Day convolution, a morphism
\[
 H(b,c)\otimes_k^{\mathbf L}H(a,b)\longrightarrow H(a,c).
\]
Here the tensor product on product-preserving presheaves is the reflected Day tensor. The reflection is monoidal: its generating localization maps are $y(U)\amalg y(V)\to y(U\oplus V)$ and $\varnothing\to y(0)$; Day tensoring with a representable $y(T)$ carries them to maps of exactly the same form, because $T\otimes_k-$ preserves finite direct sums. Tensoring with an arbitrary presheaf follows by colimits. Thus the pairing has a uniquely induced morphism after this reflection, coherently for all iterated compositions. The reflected tensor agrees with derived tensor because it sends representables associated with $U,V$ to the representable associated with $U\otimes_kV$ and preserves sifted colimits separately. Every simplicial $k$-module is a sifted colimit of finite free modules, as follows from its usual simplicial free resolution and expressing free modules as filtered colimits of their finite free submodules. These two properties determine the tensor product and the displayed morphism. Equivalently, the matrix operations show directly that the pairing is linear separately in both variables, which is precisely the relation imposed by the reflected Day tensor.

Associativity is the action associativity together with associativity of composition in $\mathcal L$; all higher compatibility diagrams are images of those same module-action diagrams. The identity of $a$ gives the unit $k\to H(a,a)$, and its unit diagrams are the action unit diagrams. This proves a coherent enriched structure. The canonical composition-compatible hammock comparison identifies the underlying spaces of these hom-objects with the nerves of corrected weak-vertical path categories. Transporting the enriched structure along that specified comparison gives the claimed refinements; no equivalence for arbitrary vertical arrows is used.
\end{proof}

\begin{result}{E.8}{Proved}{The corrected pointed refinement after localization}
Let $C$ be a small category with a zero object $0$, and let $W$ be any wide subcategory. Its infinity-localization is pointed. Its localization mapping spaces, and therefore the corrected weak-vertical path spaces with the zero-arrow basepoints, admit associative unital composition through the smash product, coherently in all variables. This assertion concerns the localized hom-spaces; it does not assert a strictly absorbing operation on every prescribed finite path stage.
\end{result}
\begin{proof}
The two functors $\{0\}\to C$ and $C\to\{0\}$ are relative, and the former is both left and right adjoint to the latter. The adjunction descent argument in E.7 therefore makes the image of $0$ both initial and terminal in $\mathcal L=C[W^{-1}]$. In particular the spaces $\operatorname{Map}_{\mathcal L}(0,a)$ and $\operatorname{Map}_{\mathcal L}(a,0)$ are contractible for every $a$.

The zero map $a\to b$ is the composite through $0$. On the wedge
\[
 (\operatorname{Map}(b,c)\times\{0\})\ \cup\
 (\{0\}\times\operatorname{Map}(a,b)),
\]
composition factors through $0$. Its contraction to the zero map is induced by the contractible incoming and outgoing mapping spaces of $0$. On the intersection the two contractions agree coherently, since both are the same composite through the zero object. The universal property of the homotopy pushout defining smash therefore gives
\[
 \operatorname{Map}(b,c)\wedge\operatorname{Map}(a,b)
 \longrightarrow\operatorname{Map}(a,c).
\]
For any finite composable sequence, the same construction contracts the entire fat wedge: if several inputs are zero, collapse all intervening maps through $0$. The spaces of choices made through $0$, and their compatibility extensions made through $0$, are contractible, because the relevant mapping spaces to and from $0$ are contractible. This is a statement about these factorizations through the zero object; it does not assert contractibility of the space of all nullhomotopies in an arbitrary mapping space. Associativity of composition supplies the face compatibilities, and identities supply degeneracies. This gives all the coherences of an enrichment over pointed spaces with smash tensor and unit $S^0$. The hammock comparison sends the one-edge zero-arrow path to the specified zero map and respects concatenation, so it identifies the corrected path hom-spaces with this pointed enrichment.
\end{proof}

The original all-vertical comparison and the literal ind-stage assertion remain separate. In particular, a weak equivalence from the colimit of based path stages to a pointed localization hom-space does not provide a uniform later stage for every required absorption homotopy on an entire infinite stage. E.8 does not supply that missing uniformity and is not recorded as a proof of the full ind-pointed assertion in PATHS C.10.
